\chapter{Methodology}
\label{ch:methodology}

This chapter specifies the experiments whose results appear in later versions of
this paper.
It is stated in advance and in full so that the design may be
evaluated independently of the outcome.

Three instruments are used, answering three different kinds of question.
Litmus testing on real hardware establishes what commercially available machines
actually do.
A configurable simulator establishes what each memory model costs,
holding everything else fixed. An axiomatic checker establishes that the simulator
implements the models it claims to implement.

\section{Instrument 1: Litmus Testing on Hardware}
\label{sec:litmus-method}

\subsection{Test Harness}

Each litmus test is realised as a C program with one pinned thread per core,
executing the test body in a loop of $N = 10^7$ iterations.
Each iteration
resets the shared variables, releases the threads, executes the body, and records
the observed final state.
The design constraints are as follows.

\paragraph{Thread placement.} Threads are pinned with
\code{pthread\_setaffinity\_np}. Placement is a controlled variable, not an
incidental one: two threads on sibling SMT contexts of the same physical core
share a store buffer and will produce different results from two threads on
distinct cores.
The topology reported by \code{lscpu -e} is recorded and reported
with every result.

\paragraph{Synchronisation between iterations.} A conventional
\code{pthread\_barrier} is unsuitable: barriers contain full fences, which suppress
precisely the window under measurement.
A lightweight sense-reversing spin barrier
built from relaxed atomics is used instead, and its own fence content is documented.

\paragraph{Preventing compiler interference.} Shared variables are declared
\code{\_Atomic} and accessed with \code{memory\_order\_relaxed}, which forbids the
compiler from eliminating or reordering the accesses while emitting no hardware
fences.
Using \code{volatile} instead is common practice but is not a substitute:
\code{volatile} constrains the compiler only with respect to other volatile
accesses.
Both variants are compiled and compared, and the difference is itself
reported.

\paragraph{Variable placement.} Each shared variable is placed on its own cache
line by padding to $64$ bytes.
Without this, false sharing changes the coherence
behaviour and therefore the observation rate.

\subsection{Platforms}

\begin{table}[H]
\centering
\small
\caption{Target platforms. Exact models to be recorded in the results chapter.}
\label{tab:platforms}
\begin{tabular}{@{}llll@{}}
\toprule
\textbf{Class} & \textbf{Model} & \textbf{Expected} & \textbf{Role} \\
\midrule
x86-64  & TSO      & SB observable; MP, LB not & strong-model baseline \\
AArch64 & ARMv8    & SB, MP, LB observable; IRIW not & weak-model comparison \\
\bottomrule
\end{tabular}
\end{table}

The AArch64 platform may be a single-board computer, an Android device under
Termux, or an institutional server; whichever is used, its microarchitecture and
core topology are recorded. The IRIW test requires at least four cores and is
reported only where the platform supports it.

\subsection{Measurements Reported}

\begin{enumerate}
\item \textbf{Observation frequency} per test per platform, over $10^7$
  iterations, with ten repetitions reported as mean and standard deviation.
  Reporting frequency rather than a binary observed/not-observed classification is
  a deliberate choice: the gap between what an architecture permits and what a
  particular implementation exhibits is one of the paper's subjects, and it is
  invisible in a binary encoding.
\item \textbf{The effect of fences.} Each test is re-run with the appropriate
  fence inserted. The expected result is that every observation frequency falls to
  zero; any test for which it does not indicates either a misplaced fence or a
  misunderstanding of the fence's semantics, and would require investigation.
\item \textbf{Fence cost in cycles}, measured by timing a tight loop with and
  without the barrier instruction using \code{rdtsc} on x86 and the cycle counter
  on AArch64, with frequency scaling disabled.
\item \textbf{Sensitivity} to thread placement and to inserted delay between the
  store and the load, establishing the width of the vulnerability window.
\end{enumerate}

\subsection{Threats to Validity}

Litmus testing establishes that a behaviour \emph{can} occur; it can never
establish that one cannot.
A zero observation count is consistent with the
behaviour being architecturally forbidden, with it being permitted but not
exhibited by this implementation, and with the harness suppressing it.
This
paper therefore reports observations as lower bounds and pairs every measurement
with the corresponding architectural statement from the vendor manual. Where the
two disagree --- permitted but never observed --- the discrepancy is reported as a
finding rather than resolved by assumption.

\section{Instrument 2: The Simulator}
\label{sec:sim-method}

\subsection{Rationale}

Hardware answers what x86 and ARM do.
It cannot answer what SC would cost, because
no SC hardware exists to measure.
That question requires a machine in which the
memory model is a parameter.

\subsection{Architecture}

The simulator models $N$ cores, each with a private L1, connected by a directory
protocol over an interconnect of configurable latency.
Its distinguishing feature
is that the ordering policy is an interchangeable component: a single interface
decides, for each memory operation, whether it may issue given the operations
still outstanding.
Four implementations of that interface are provided --- SC-naive,
SC-speculative, TSO and Weak --- and \emph{no other part of the simulator differs
between configurations}.
The memory model is thereby the only independent variable.

\begin{table}[H]
\centering
\small
\caption{Simulator configuration parameters and their default values.}
\label{tab:simconfig}
\begin{tabular}{@{}lll@{}}
\toprule
\textbf{Parameter} & \textbf{Default} & \textbf{Swept over} \\
\midrule
Cores                & 8            & 2, 4, 8, 16, 32 \\
Store buffer entries & 32           & 8, 16, 32, 64 \\
Load queue entries   & 32           & 16, 32, 64 \\
L1 size / assoc.     & 32\,kB / 8   & fixed \\
Coherence            & MESI directory & fixed \\
Interconnect latency & 20 cycles    & 10, 20, 40 \\
Memory latency       & 200 cycles   & fixed \\
Memory model         & ---          & SC-naive, SC-spec, TSO, Weak \\
\bottomrule
\end{tabular}
\end{table}

\subsection{Speculative SC}

The SC-speculative policy is the one whose viability the paper tests.
Its
operation is as follows.

\begin{enumerate}
\item A load is permitted to issue even where SC would require it to wait, and its
  address is retained in the load queue marked speculative.
\item Incoming coherence invalidations are matched against speculative load queue
  entries.
\item A match indicates that a remote store to a speculatively read line occurred
  within the window, so the speculation was potentially observable.
The load and
  all younger instructions are squashed and re-executed.
\item Absent a match, the load retires and the relaxation is unobservable.
\end{enumerate}

Both required mechanisms already exist in the modelled core: coherence
invalidations are generated by the protocol regardless, and the squash path is the
branch-misprediction recovery path.
The added hardware is an address comparator
against the load queue.

\subsection{Workloads}

Standard benchmark suites are unsuitable as the primary workload.
SPEC and PARSEC
are dominated by cache and branch behaviour; the memory model contributes little
and the effect under study would be buried in noise.
Workloads are therefore
selected for communication density.

\begin{table}[H]
\centering
\small
\caption{Simulated workloads.}
\label{tab:workloads}
\begin{tabular}{@{}llp{6.4cm}@{}}
\toprule
\textbf{Tier} & \textbf{Workload} & \textbf{Stresses} \\
\midrule
Sync    & Spinlock counter        & Fence cost on acquire and release under contention \\
Sync    & Dekker mutual exclusion & Store~$\to$~load ordering directly \\
Sync    & SPSC ring buffer        & Acquire/release pairs in steady state \\
Sync    & Treiber stack           & Atomic RMW and ordering around it \\
Sync    & Sense-reversing barrier & All-core synchronisation \\
Kernel  & Parallel reduction      & Mostly-private access with periodic combining \\
Kernel  & Shared-frontier BFS     & Irregular sharing \\
Kernel  & Atomic histogram        & Contended atomic updates \\
\bottomrule
\end{tabular}
\end{table}

The two tiers serve different purposes.
The synchronisation microbenchmarks
isolate the effect; the kernels test whether it survives into code that performs
useful work.
It is anticipated that the gap between models narrows substantially
in the second tier, and that outcome, if observed, is a result and not a
disappointment.

\subsection{Metrics}

Per run, the simulator emits: total cycles; IPC; stall cycles decomposed by cause
(store buffer full, fence drain, cache miss, coherence wait); squash count;
coherence messages; mean store buffer occupancy.

The headline metric is execution time normalised to Weak.
The explanatory metric
is the decomposition of stall cycles: the hypothesis under test predicts that
SC-naive's ordering stalls are largely \emph{converted} into squash cycles under
SC-speculative, and that squash cycles are fewer because they are incurred only on
genuine conflicts.
A result without this decomposition would establish a
correlation without a mechanism.

\subsection{Experimental Protocol}

Each configuration is run once as warm-up and discarded, then ten times with
distinct seeds. Results are reported as mean and standard deviation with error
bars on all figures.
The full sweep is $8$ workloads $\times$ $4$ models $\times$
$5$ core counts $\times$ $10$ repetitions, or $1600$ runs, driven by a single
script and emitting one CSV row per run.
All figures are generated from those CSVs
by scripts under version control, so that the figures are reproducible from the
data by one command.

\paragraph{A methodological commitment.} Every number in the results chapters
originates from this simulator under the configuration recorded alongside it.
No
comparison is made against published figures obtained on other simulators or other
hardware, since differing microarchitectural assumptions would render such
comparisons meaningless.
The value of the instrument lies entirely in its being a
controlled experiment.

\section{Instrument 3: Verification}
\label{sec:verif-method}

A performance comparison between memory models is worthless if the implementations
do not in fact implement those models.
The third instrument establishes that they
do.

For each litmus test and each model, two sets of final states are computed:

\begin{enumerate}
\item The \emph{permitted} set, by enumerating all candidate executions and
  applying the acyclicity criterion of \cref{def:permitted} directly.
\item The \emph{produced} set, by exhaustively enumerating all interleavings the
  simulator admits for that test under that ordering policy.
\end{enumerate}

The two sets must be equal. A produced state outside the permitted set is a bug in
the ordering policy.
A permitted state never produced indicates that the
implementation is stricter than the model requires --- not a correctness error, but
a fact that must be reported, since it would bias the performance comparison in
that model's favour.

Litmus tests are small enough that exhaustive enumeration is tractable, and the
checker is a distinct piece of software from the timing simulator, so that a
common bug is unlikely to mask itself in both.

\section{Tooling and Reproducibility}

The simulator is written in C++ and emits CSV. Figures are produced in Python with
pandas and matplotlib, one script per figure, exported as vector PDF. Litmus tests
are C11 with pthreads.
The paper is written in \LaTeX{} with all figures in TikZ,
and references are maintained in Zotero and exported to BibTeX. Random seeds are
fixed, a single \code{run\_all.sh} reproduces every result, and the repository is
cited in the final version.

\section{Anticipated Results and Falsifiability}

The hypothesis under test is that speculative sequential consistency is
substantially cheaper than naive sequential consistency, and close to weak
ordering, on a core that already speculates.
Concretely, the paper predicts:

\begin{itemize}
\item SC-naive between $1.5\times$ and $2\times$ the execution time of Weak on the
  synchronisation tier;
\item SC-speculative within $1.05\times$ to $1.15\times$ of Weak on the same tier;
\item the advantage of SC-speculative narrowing as core count rises, driven by
  squash rate.
\end{itemize}

Stating these in advance makes the study falsifiable.
Should SC-speculative prove
expensive at all core counts, or should its cost prove insensitive to squash rate,
the hypothesis is refuted and the historical verdict is vindicated --- which would
be a result of equal interest and is reported as such.
